\documentclass[11pt,a4paper]{article}
\usepackage[utf8]{inputenc}\usepackage[T1]{fontenc}
\usepackage[a4paper,margin=2.5cm]{geometry}
\usepackage{microtype,hyperref,xcolor}
\hypersetup{colorlinks=true,linkcolor=blue,urlcolor=blue}
\setlength{\parskip}{0.4em}\setlength{\parindent}{0pt}\sloppy
\begin{document}

\textbf{Cover letter to the Editors of \emph{Quantum Science and Technology}}

\vspace{0.3cm}
Hikaru Wakaura, Taiki Tanimae\\
QIRI (Quantum Integrated Research Institute Inc.), 1-16-3, Akasaka, Minato-ku, Tokyo 107-0061, Japan\\
\href{mailto:h.wakaura@qiri.co.jp}{h.wakaura@qiri.co.jp} \quad ORCID: 0009-0001-2341-7710

\vspace{0.2cm}
September 14, 2026

\vspace{0.3cm}
Dear Editors,

On September 14, 2026, we submitted our manuscript ``\textbf{Blind catalytic quantum error correction: target-state estimation and fidelity bounds under an idealized catalytic recovery interface}'', authored by Hikaru Wakaura and Taiki Tanimae, for consideration as a Paper in \emph{Quantum Science and Technology}.
The submission consists of the main manuscript (17 figures, 9 tables) and a frozen release of the open-source package \texttt{blind\_cqec} v0.2.0 (MIT license, \url{https://github.com/deeptell-inc/blind_cqec_pkg}) with a 91-test reproducibility suite, shipped raw results files, and a script-to-figure map.

\textbf{What we discovered.}
Catalytic quantum error correction (CQEC) is a threshold-free state-recovery framework whose original formulation requires the ideal target state to be known exactly.
In this study we show that the target-knowledge requirement can be displaced onto a classical estimation problem.
Under an explicitly stated idealized recovery interface---the mode-inclusion-restricted projection of the estimate---recovery fidelity is bounded by estimation fidelity plus a mode-survival penalty, and an 84-condition benchmark realizes both regimes of that bound.
The resulting estimator design-space map is sharp: with exactly calibrated noise, channel inversion coincides with the oracle at every dimension and both are capped by a mode-survival ceiling; the noise-model-free alternative, coherence maximization, matches the ceiling to within $1.5\%$ at $d \leq 16$ and degrades with dimension; and the operational choice between the two is decided by calibration accuracy, with the dephasing rate exponentially dominant.

\textbf{Why this matters.}
The result separates, for the CQEC setting, what is a classical estimation question from what is a structural limit of catalytic recovery, and it reports the negative results with the same prominence as the positive ones: coherence maximization is actively harmful for mixed targets, no strategy reaches high fidelity where mode annihilation sets the ceiling, and universal single-copy recovery is impossible for any deterministic target-independent apparatus.
This map is the prerequisite for the two open problems the manuscript states as future work---an explicit joint covariant recovery channel and an operational measurement model---and it is directly useful to practitioners choosing an estimator for variational and iterative algorithms whose ideal output is unknown to the correction module.

\textbf{Why \emph{Quantum Science and Technology}.}
The manuscript sits at the interface of quantum resource theory, quantum state estimation, and near-term error handling, which is the journal's core scope, and it is accompanied by a fully reproducible open-source implementation with raw data and a figure-by-figure reproduction map, matching the journal's emphasis on practically relevant, verifiable results.
The claims are deliberately scoped to what is demonstrated, and the scope is stated in the abstract, methods, and limitations.

\textbf{Suggested referees.}
\begin{enumerate}\setlength{\itemsep}{0pt}
\item Ryuji Takagi, Nanyang Technological University, Singapore (\texttt{ryuji.takagi@ntu.edu.sg}) --- catalytic coherence transformations.
\item Mark M.\ Wilde, Cornell University, USA (\texttt{wilde@cornell.edu}) --- quantum Shannon theory, channel fidelity bounds.
\item Andreas Winter, Universitat Aut\`onoma de Barcelona, Spain (\texttt{andreas.winter@uab.cat}) --- catalytic resource theories.
\item B\'alint Koczor, University of Oxford, UK (\texttt{balint.koczor@maths.ox.ac.uk}) --- error mitigation, near-term VQE.
\item Daniel Stilck Fran\c{c}a, ENS Lyon / CNRS, France (\texttt{daniel.stilck-franca@ens-lyon.fr}) --- quantum state estimation, copy complexity.
\item Iman Marvian, Duke University, USA (\texttt{iman.marvian@duke.edu}) --- asymmetry resource theory, covariant operations.
\end{enumerate}
\textbf{Non-preferred referees.} None.

\textbf{Statements.}
This manuscript is not under consideration elsewhere. A previous version was reviewed at another journal; the present version incorporates the resulting corrections, in particular the explicit statement that the catalytic transformation itself is not simulated and that all resource counts are conditional on oracle density-matrix access.
Both authors have read and approved the submission. We declare no competing interests.

\vspace{0.4cm}
Sincerely,\\[0.4em]
Hikaru Wakaura, Taiki Tanimae\\
QIRI Inc., Tokyo
\end{document}
